\chapter{Fundamentos de los métodos metaheurísticos}
\label{ch:fundamentos-metaheuristicos}

\section{Introducción del capítulo}
Las metaheurísticas surgieron como respuesta a problemas de optimización donde los métodos
exactos, deterministas o basados en gradiente no resultaban prácticos. En esos contextos, lo
importante no es garantizar una solución cerrada, sino encontrar soluciones de alta calidad en
espacios amplios, no lineales y frecuentemente multimodales. La sintonización PID en
microredes aisladas encaja bien en esta categoría, ya que combina restricciones físicas,
objetivos dinámicos y fuerte interacción entre parámetros.

\section{Objetivo del capítulo}
Este capítulo presenta la lógica común que comparten los métodos metaheurísticos y ofrece el
vocabulario mínimo necesario para estudiar después algoritmos particulares como PSO, GWO,
Eagle Strategy y ABC. Se busca mostrar que, más allá de su inspiración biológica o social,
todos persiguen un equilibrio entre explorar alternativas nuevas y explotar las soluciones que
ya demostraron ser prometedoras.

\section{Concepto de metaheurística}
Una metaheurística puede definirse como una estrategia general de búsqueda capaz de guiar
la exploración de soluciones en problemas complejos sin requerir información exacta sobre la
forma analítica del óptimo. En lugar de seguir una trayectoria única, estas técnicas utilizan
poblaciones, reglas estocásticas, memoria parcial y criterios de actualización que permiten
recorrer el espacio de búsqueda con flexibilidad. Por ello resultan especialmente útiles en
aplicaciones donde la función objetivo depende de simulaciones o evaluaciones costosas.

\section{Exploración y explotación}
El corazón de toda metaheurística es el equilibrio entre exploración y explotación. Explorar
significa visitar regiones nuevas del espacio de búsqueda para evitar quedar atrapado en
óptimos locales; explotar significa refinar las mejores soluciones conocidas para mejorar su
calidad. Si el algoritmo explora demasiado, tarda en converger; si explota demasiado pronto,
puede estancarse en soluciones mediocres. El diseño de parámetros y operadores busca
justamente administrar esa tensión.

\section{Ventajas frente a métodos clásicos}
Frente a reglas clásicas como Ziegler--Nichols, las metaheurísticas ofrecen la ventaja de buscar
ganancias directamente sobre el desempeño observado en simulación o experimento. Esto les
permite incorporar restricciones, múltiples métricas y comportamientos no lineales sin reducir
la planta a una representación excesivamente simplificada. Aunque requieren más cómputo, ese
costo suele ser aceptable cuando se busca una sintonización fuera de línea que luego será
implementada en un controlador industrial convencional.

\section{Clasificación general}
De manera general, las metaheurísticas pueden agruparse en métodos evolutivos, métodos de
enjambre e híbridos. Los primeros se inspiran en selección, cruce y mutación; los segundos
modelan cooperación social, colonias o patrones colectivos; los híbridos combinan operadores
de diferentes familias para mejorar convergencia o robustez. En esta obra se prioriza la familia
de enjambre porque ha mostrado buenos resultados en control de frecuencia y porque permite
una implementación comparativa relativamente homogénea entre algoritmos.

\section{Cierre del capítulo}
Con esta base conceptual, ya es posible estudiar algoritmos concretos sin perder de vista la
idea general que los conecta. El capítulo siguiente desarrolla el método PSO, que sirve como
punto de referencia principal en la propuesta comparativa del libro.
